\begin{notechapter}
\chapter{Gradient, Hessian and Taylor's Formula}
\label{ch:gradient-hessian}
\chaptersubtitle{%
  What an inner product adds, and the second-order structure}

\begin{chapterplan}
\begin{itemize}[nosep]
  \item The gradient as the Riesz representative of the differential; its
        dependence on the inner product; steepest ascent; orthogonality to level
        sets.
  \item Gradient descent passes the origin test but depends on the metric.
  \item Higher differentials; symmetry of the second differential; the Hessian
        as a bilinear form (intrinsic) and as an operator (needs a metric).
  \item Taylor's formula with remainder; critical points and second-order
        conditions.
  \item Convexity, strong convexity and $L$-smoothness; the descent lemma;
        convergence rates of gradient descent.
  \item Newton's method is invariant under affine reparametrizations; gradient
        descent is not.
\end{itemize}
\smallskip\noindent\textbf{Machine learning threads.}
\begin{itemize}[nosep]
  \item Bounds on the learning rate from smoothness.
  \item Hessian--vector products.
  \item Preconditioners, from diagonal scalings to the natural gradient.
\end{itemize}
\end{chapterplan}
\end{notechapter}
